\section{Economics and Mathematics behind the FEDWatcher}

\subsection{Theory}
The object of interest of this project is the outcome of the next FOMC meeting, defined as the change in the upper bound of the federal funds target range relative to the current one. In order to compute the prediction, we described the possible actions of the FED through five buckets:
\[
  Y_{t+1} \in \mathcal{Y} = \{-50,\,-25,\,0,\,+25,\,+50\}\ \text{basis points}.
\]
Then, the estimation model is a Probit that follows this specification:
\[
  Y^*_t = X_t \beta + \epsilon_t,
\]
and the conditioning set is
\[
  \mathcal{X}_t = \underbrace{\bigl(\tau_t^{(1)}, \ldots, \tau_t^{(K)}\bigr)}_{\text{LLM tone scores}} \cup \underbrace{\bigl(\pi_t,\, \pi_t^{\text{core}},\, u_t,\, u_t - u_t^\ast,\, y_t^{2y},\, r_t^{\text{ff}},\, y_t^{2y} - r_t^{\text{ff}}\bigr)}_{\text{macro and rates}}.
\]
Here $\pi_t$ and $\pi_t^{\text{core}}$ are year-on-year headline and core CPI, $u_t$ is unemployment, $u_t^\ast$ is the natural rate (FRED \texttt{NROU}) so that $u_t - u_t^\ast$ is the unemployment gap, $y_t^{2y}$ is the 2-year Treasury yield, $r_t^{\text{ff}}$ is the current Fed-funds target, and $y_t^{2y} - r_t^{\text{ff}}$ is a crude proxy for the market's pricing of forthcoming policy.

\subsection{Tone Extraction}
For each new document $d$, AnalystAgent splits the text into $K\!\approx\!5$ sections using regex anchors trained on the structure of FOMC statements (assessment, outlook, policy stance, voting record). The LLM is then prompted with a strict JSON schema and returns, per section, an overall tone in $[-1,1]$, an inflation assessment label (\texttt{cooling}, \texttt{stable}, \texttt{heating}), a labor-market assessment, a forward-guidance label (\texttt{dovish}, \texttt{neutral}, \texttt{hawkish}) and a confidence score. Section scores are aggregated to a single document tone by a calibrated weighted average:
\[
  \tau_t = \frac{\sum_k w_k \, \text{conf}_k \, s_k}{\sum_k w_k \, \text{conf}_k}, \qquad w_k \text{ from } \texttt{section\_weight\_calibration},
\]
where the section weights $w_k$ are estimated by regressing historical document tone against the realized next-meeting outcome on the training window. The dual-model analyst runs two distinct LLMs in parallel and takes a confidence-weighted average to reduce single-model variance; the resulting tone is stored in \texttt{sentiment2}/\texttt{sentiment3}, and a weighted variant in \texttt{sentiment\_w}.

\subsection{Nowcast Model}
The baseline nowcast is an \emph{ordered probit} over \emph{j} with cutpoints $c_1 < c_2 < \cdots < c_5$:
\[
  Y^*_t = X_t \beta + \epsilon_t,\qquad \Pr(Y_{t} = j \mid \mathcal{X}_t) = \frac{\exp{\mathcal{X}_t}\beta_j}{1+\sum_{k=1}^J\exp{\mathcal{X}_t\beta_j}}),\qquad c_0 = -\infty,\ c_6 = +\infty,
\]
The full specification is then defined as
\[
  Y^*_t = \beta_1 \cdot S_t + \beta_2 \cdot \Delta CPI_t + \beta_3 \cdot Ugap_t + \epsilon_t
\]
where $\Delta CPI_t$ is the variation in the Consumer Price Index, $Ugap_t$ is the unemployment gap and $\epsilon_t$ is the error term. $S_t$ represents the tone score, extracted by a text analysis performed by an LLM from all the statements and documents published by the FED. In particular, $S_t$ is an Exponentially Weighted Moving Average (EWMA) defined as follows:

\[
  \text{S}_t = \alpha_t \cdot \text{tone\_score}_t + (1 - \alpha_t) \cdot \text{S}_{t-1},\quad \alpha_t = 1 - e^{-\lambda \Delta t}, \quad \text{where } \lambda = \frac{\ln(2)}{h}.
\]
The first observation in our database seeds the series with alpha = 1.0 and $S_{t-1}$ =  $S_t$. After some research and literature review, we decided to set \emph{h} equal to 21. We believe this is the best option for our purposes the project objectives.

\subsection{Computation of Tone Implied Rate and Signal Divergence}
The final tone implied rate $\widehat{r}^{\text{tone}}_{t+1}$ is computed by summing to the current Fed-funds target the sum of the probabilities of each bucket $j$ multiplied by the corresponding basis points:

\[
  \widehat{r}^{\text{tone}}_{t+1} = r_t^{\text{ff}} + \sum_{j} \Pr(Y_{t} = j \mid \mathcal{X}_t) \cdot \text{bp}_j.
\]
The StrategistAgent then computes the divergence
\[
  \delta_t = r^{\text{mkt}}_{t+1} - \widehat{r}^{\text{tone}}_{t+1}\,
\]
where $r^{\text{mkt}}_{t+1}$ is the market-implied next-meeting rate from Fed-funds futures or, when unavailable, from the 1m OIS and $\widehat{r}^{\text{tone}}_{t+1}$ is the tone-implied next-meeting rate (a calibrated mapping from $\tau_t$ to basis points).

\section{Data Science}

\subsection{Data Sources}
All macro series come from FRED via the public REST API: \texttt{DFEDTARU}, \texttt{DFEDTARL}, \texttt{DFF}, \texttt{FEDFUNDS}, \texttt{CPIAUCSL}, \texttt{CPILFESL}, \texttt{UNRATE}, \texttt{NROU}, \texttt{DGS2}, \texttt{SOFR}. FOMC statements are scraped from \url{federalreserve.gov} and from the synthetic \texttt{fakefed.ellep.it} used for parser tests; minutes and speeches are excluded from ingestion. The historical statement backfill covers 1994--present.

\subsection{Why SQLite}
The schema (\texttt{db/schema.sql}) deliberately favors clarity over performance: every \texttt{sentiment*} table uses \texttt{document\_id} as a foreign key to \texttt{documents}, and signals are reconstructable joins rather than materialized views. SQLite was chosen over Postgres because the entire dataset fits in $\sim$4\,MB, the read pattern is dashboard-style (few writers, many readers), and a single file is trivially backupable and reproducible.
